\subsection{Diffusions}

DEFINITION 3. --- Let X be a locally compact space, and let \(\Lambda : t \mapsto \lambda_t\) be a mapping of T into \(\mathcal{M}_{+}(X)\) . The mapping \(\Lambda\) is said to be a diffusion of T in X if \(\Lambda\) is adequate for every positive measure on T with compact support. The diffusion \(\Lambda\) is said to be bounded if all of the measures \(\lambda_{t}\) are bounded and \(\sup_{t\in T}\|\lambda_{t}\|<+\infty\) ; this quantity is then called the norm of \(\Lambda\) and is denoted \(\|\Lambda\|\) .

The following proposition merely translates the definition:

PROPOSITION 8. --- For a mapping \(\Lambda: t \mapsto \lambda_t\) of T into \(\mathcal{M}_+(\mathrm{X})\) to be a diffusion, it is necessary and sufficient that the following conditions be satisfied:

\begin{enumerate}
\def\labelenumi{\arabic{enumi})}
\item
  For every lower semi-continuous function \(f \geqslant 0\) defined on \(X\) , the function \(t \mapsto \lambda_t^\bullet(f)\) is universally measurable on \(T\) .
\item
  For every function \(g \in \mathcal{K}_{+}(X)\) , the function \(t \mapsto \lambda_t(g)\) is locally bounded in \(T\) .
\item
  For every lower semi-continuous function \(f \geqslant 0\) defined on X and for every positive measure \(\mu\) on T with compact support, the following relation holds, where \(\nu\) denotes \(\int \lambda_{t} d\mu(t)\) :
\end{enumerate}

\[
\int^ {\bullet} f (x) d \nu (x) = \int^ {\bullet} d \mu (t) \int^ {\bullet} f (x) d \lambda_ {t} (x).\tag{12}
\]

Suppose that \(\Lambda\) is a diffusion. The condition 1) is then satisfied by the definition of adequate mappings (No.~1, Def. 1) and Prop. 6; the condition 3) is satisfied by formula (4), since \(\Lambda\) is \(\mu\) -adequate. Let \(g \in \mathcal{K}_{+}(X)\) and let \(u\) be the function \(t \mapsto \lambda_t(g)\) (universally measurable, by 1)); suppose that \(u\) is not locally bounded. There would then exist a compact set \(K\) such that \(u\) is not bounded on \(K\) , hence there would exist a sequence \((t_n)\) of elements of \(K\) such that \(u(t_n) \geqslant n^2\) for all \(n \geqslant 1\) ; then \(u\) is not integrable for the measure \(\mu = \sum_{n \geqslant 1} \frac{1}{n^2} \varepsilon_{t_n}\) with compact support, contrary to the hypothesis on \(\Lambda\) , which implies that \(t \mapsto \lambda_t(g)\) is integrable for every positive measure with compact support. The above three conditions are thus necessary. Conversely, the conditions 1) and 2) imply that \(\Lambda\) is scalarly essentially \(\mu\) -integrable for every measure \(\mu\) with compact support. Since every measure \(\mu' \geqslant 0\) that is bounded above by a measure \(\mu\) with compact support also has compact support, the conditions 1) and 3) express that \(\Lambda\) is \(\mu\) -adequate for every positive measure with compact support, which is indeed the sought-for result.

PROPOSITION 9. --- Let \(\Lambda : t \mapsto \lambda_t\) be a mapping of T into \(\mathcal{M}_+(\mathrm{X})\) , such that the function \(t \mapsto \lambda_t(g)\) is universally measurable and locally bounded in T for every \(g \in \mathcal{K}_+(\mathrm{X})\) . One can affirm that \(\Lambda\) is a diffusion in each of the following cases:

\begin{enumerate}
\def\labelenumi{\alph{enumi})}
\item
  the topology of X admits a countable base;
\item
  \(\Lambda\) is universally measurable for the vague topology.
\end{enumerate}

For, let \(\mu\) be a positive measure on T with compact support; the mapping \(\Lambda\) is scalarly essentially \(\mu\) -integrable, hence \(\mu\) -adequate if either a) or b) is satisfied (No.~1, Prop. 2).

For the rest of this section, we shall adopt the following notations: we will denote by \(\langle\eta,h\rangle\) the upper essential integral, for a positive measure \(\eta\) , of a positive \(\eta\) -measurable function h. The mapping \(\Lambda:t\mapsto\lambda_{t}\) will be a diffusion of T in X. If f is a positive universally measurable function defined on X, we shall denote by \(\Lambda f\) the mapping \(t\mapsto\lambda_{t}^{\bullet}(f)\) . If \(\mu\) is a positive measure on T such that \(\Lambda\) is scalarly essentially \(\mu\) -integrable, we shall denote by \(\mu\Lambda\) the measure \(\int\lambda_{t}d\mu(t)\) . The definition of the integral then takes the form

\[
\langle \mu \Lambda , f \rangle = \langle \mu , \Lambda f \rangle \quad \text { for } f \in \mathscr {K} _ {+} (\mathrm{X}).
\]

We shall say that a positive measure \(\mu\) on T belongs to the domain of \(\Lambda\) if \(\Lambda\) is \(\mu\) -adequate: this amounts to saying (in view of Prop. 8) that \(\Lambda\) is scalarly essentially \(\mu\) -integrable and \(\langle\mu'\Lambda,f\rangle=\langle\mu',\Lambda f\rangle\) for every positive measure \(\mu'\leqslant\mu\) and every lower semi-continuous positive function f.

PROPOSITION 10. --- Let f, g be two positive universally measurable functions on X, let a be a number \(\geqslant 0\) , and let \(\mu\) and \(\nu\) be two positive measures on T. Then:

\begin{enumerate}
\def\labelenumi{\alph{enumi})}
\item
  \(\Lambda (f + g) = \Lambda f + \Lambda g, \Lambda (af) = a\Lambda f.\)
\item
  If \(\mu\) and \(\nu\) belong to the domain of \(\Lambda\) , then so do \(\mu + \nu\) and \(a\mu\) , and one has \((\mu + \nu)\Lambda = \mu \Lambda + \nu \Lambda\) , \((a\mu)\Lambda = a(\mu \Lambda)\) .
\end{enumerate}

The only non-obvious point is that \(\mu + \nu\) belongs to the domain of \(\Lambda\) , which is treated by observing that every positive measure bounded above by \(\mu + \nu\) is of the form \(\mu' + \nu'\) , where \(\mu' \leqslant \mu\) , \(\nu' \leqslant \nu\) (the `decomposition lemma', Ch. II, §1, No.~1). See also the next proposition.

PROPOSITION 11. --- For a positive measure \(\mu\) on \(T\) to belong to the domain of \(\Lambda\) , it is necessary and sufficient that \(\Lambda\) be scalarly essentially \(\mu\) -integrable.

This condition is obviously necessary. Conversely, suppose it is satisfied, and let f be a lower semi-continuous positive function defined on X. The function \(\Lambda f\) is universally measurable, hence \(\mu\) -measurable. We are going to prove that \(\langle\mu,\Lambda f\rangle=\langle\mu\Lambda,f\rangle\) ; since this equality will also be valid for every positive measure \(\mu'\leqslant\mu\) , because \(\Lambda\) is also scalarly essentially \(\mu'\) -integrable, it will follow that \(\Lambda\) is \(\mu\) -adequate.

Let \((\mu_i)_{i\in I}\) be a summable family of positive measures with compact support, such that \(\mu = \sum_{i\in I}\mu_i\) (§2, No.~3, Prop. 4); the family of measures \(\mu_{i}\Lambda\) is then summable, and \(\mu \Lambda = \sum_{i\in I}\mu_{i}\Lambda\) (No.~1, Cor. of Prop. 1). Consequently \(\langle \mu \Lambda ,f\rangle = \sum_{i\in I}\langle \mu_i\Lambda ,f\rangle\) (§2, No.~2, Prop. 1); but \(\Lambda\) is \(\mu_{i}\) -adequate, thus \(\langle \mu_i\Lambda ,f\rangle = \langle \mu_i,\Lambda f\rangle\) . Again applying Prop. 1 of §2, we obtain the sought-for equality:

\[
\langle \mu \Lambda , f \rangle = \sum_ {i \in I} \langle \mu_ {i} \Lambda , f \rangle = \sum_ {i \in I} \langle \mu_ {i}, \Lambda f \rangle = \langle \mu , \Lambda f \rangle .
\]

COROLLARY 1. --- If \(\Lambda\) is a bounded diffusion, then every bounded positive measure \(\mu\) belongs to the domain of \(\Lambda\) , and \(\|\mu\Lambda\| \leqslant \|\mu\|\|\Lambda\|\) .

COROLLARY 2. --- Suppose that \(\mu\) is the sum of a summable family \((\mu_{\alpha})_{\alpha \in \mathbb{A}}\) of positive measures belonging to the domain of \(\Lambda\) . For \(\mu\) to belong to the domain of \(\Lambda\) , it is necessary and sufficient that the family of measures \(\mu_{\alpha}\Lambda\) be summable, in which case \(\mu \Lambda = \sum \mu_{\alpha}\Lambda\) .

It suffices to apply the Corollary of Prop. 1 of No.~1.

Proposition 5, expressed in the language of diffusions, takes the following form:

PROPOSITION 12. --- Let \(\mu\) be a positive measure on T that belongs to the domain of \(\Lambda\) , and let \(f\) be a universally measurable function \(\geqslant 0\) defined on X. If \(f\) is moderated for the measure \(\mu\Lambda\) , or if the measures \(\lambda_t\) are bounded, then the function \(\Lambda f\) is \(\mu\) -measurable and

\[
\langle \mu \Lambda , f \rangle = \langle \mu , \Lambda f \rangle .\tag{13}
\]

COROLLARY. --- If X is countable at infinity, or if the measures \(\lambda_{t}\) are bounded, then the function \(\Lambda f\) is universally measurable on T for every universally measurable function \(f \geqslant 0\) defined on X, and (13) holds.

